\section{Exact balancing and holomorphic polystability}
\label{sec:balancing}

The finite residual is the moment map of an actual symplectic
form on the matrix parameter space. We first normalize this form
by an equivariant change of real coordinates. A separate
correction within complex orbits then identifies the holomorphic
polystable locus.

\subsection{The actual Hamiltonian structure}

On the affine space of unitary connections on $(E^{\rm sp},h_0)$,
with anti-Hermitian endomorphism-valued one-form tangents, put
\begin{equation}\label{eq:connection-symplectic}
 \Omega_{\rm conn}(a,b)
       =-\int_{X_p}\tr(a\wedge b)\wedge\frac{\omega_p^2}{2}.
\end{equation}
This is a constant closed real form. For the gauge action
$D\mapsto kDk^{-1}$ the infinitesimal vector is
$Y_\xi=-D\xi$. In the convention $\dd H_\xi=\iota_{Y_\xi}\Omega$,
its Hamiltonian is
\[
 H_\xi(\mathcal A)
       =-\int_{X_p}\tr(\xi F_{\mathcal A})\wedge\frac{\omega_p^2}{2}.
\]
Indeed differentiation in direction $a$ and integration by
parts give
$\dd H_\xi(a)=\int\tr(D\xi\wedge a)\omega_p^2/2
=\Omega_{\rm conn}(Y_\xi,a)$.
This is the usual connection moment map
\cite[Section 1.1 and Lemma 1.5]{ACG}, on the affine space
of all unitary connections.

Pull back by the actual map $w\mapsto\mathcal A_p(w)$.
The resulting closed $K$-invariant form $\Omega_p$ has,
for $S\in\mathfrak h$, Hamiltonian
\begin{equation}\label{eq:actual-paired-hamiltonian}
       H_{\mathrm iS,p}(w)=-V\tr(S\Psi_p(w)).
\end{equation}
At zero the linear connection map is
$\eta_p(v)-\eta_p(v)^*$, so the type calculation in a
unitary coframe gives
\begin{equation}\label{eq:flat-pullback-form}
 \Omega_{0,p}(v,v'):=\Omega_p(0)(v,v')
 =2\operatorname{Im}\{c_+\tr(v_A^*v_A')+c_-\tr(v_B^*v_B')\}.
\end{equation}
The sign agrees with the convention that the circle action
on $(\CC,2\dd x\wedge\dd y)$ has Hamiltonian $-|z|^2$.
The constant form \eqref{eq:flat-pullback-form} is uniformly
nondegenerate. The connection derivative estimates prove
$\|\Omega_p(w)-\Omega_{0,p}\|\le C|w|$, with uniform
higher matrix derivatives. Thus $\Omega_p$ is symplectic
on one common small ball. The quadratic Hamiltonian of
the same action for $\Omega_{0,p}$ is
$-V\tr(S\mu_{2,p}(z))$, by direct differentiation of
\eqref{eq:quadratic-moment}.

\begin{proposition}\label{prop:exact-balancing-coordinates}
There is a smooth $K$-equivariant diffeomorphism $\chi_p$
on a common neighbourhood of zero, whose image contains
a common neighbourhood, such that
\begin{equation}\label{eq:exact-balancing}
 \Psi_p(\chi_p(z))=\mu_{2,p}(z),\qquad
 \chi_p(z)=z+O(|z|^2),\quad d\chi_p(z)=I+O(|z|).
\end{equation}
Its fixed-order higher matrix derivatives are uniformly
bounded on smaller common balls. For every
$z\in\mu_{2,p}^{-1}(0)$ there is an actual HE connection
$\mathcal A_p(\chi_p(z))$, and its complete small spectrum
is given by \eqref{eq:finite-bundle-energy} and $M$ with
relative error $O(|z|)$, including its changing kernel.
\end{proposition}

\begin{proof}
Set $\sigma_p=\Omega_p-\Omega_{0,p}$. The radial homotopy
one-form
\[
 \alpha_p(z)(v)=\int_0^1 t\,\sigma_p(tz)(z,v)\,\dd t
\]
is $K$-invariant, satisfies $\dd\alpha_p=\sigma_p$, and
has size $O(|z|^2)$ and derivative $O(|z|)$.
For $0\le t\le1$ the forms
$\Omega_{p,t}=\Omega_{0,p}+t\sigma_p$ are uniformly
nondegenerate on a common small ball. Define
$\iota_{Z_{p,t}}\Omega_{p,t}=-\alpha_p$.
Then $Z_{p,t}=O(|z|^2)$, $dZ_{p,t}=O(|z|)$, and it
is equivariant. Its flow $\phi_{p,t}$ exists to time one
on a smaller common ball: the scalar bound
$r'\le Cr^2$ keeps its trajectories inside the original
ball. Differentiating the ODE gives
$\phi_{p,1}=I+O(|z|^2)$ and $d\phi_{p,1}=I+O(|z|)$,
with the stated higher bounds. The reverse flow shows
that its image contains a common neighbourhood.
Cartan's formula gives
\[
 \frac{\dd}{\dd t}\phi_{p,t}^*\Omega_{p,t}
   =\phi_{p,t}^*(\sigma_p+\dd\iota_{Z_{p,t}}\Omega_{p,t})=0.
\]
Thus $\chi_p:=\phi_{p,1}$ pulls $\Omega_p$ back to
$\Omega_{0,p}$.

Equivariance identifies the differentials of
$H_{\mathrm iS,p}\circ\chi_p$ and
$-V\tr(S\mu_{2,p})$. Both vanish at zero, so the
Hamiltonians agree exactly. The trace pairing on
$\mathfrak h$ is nondegenerate, proving
\eqref{eq:exact-balancing}. The diagonal central circle
has identically zero Hamiltonian and leaves no scalar
ambiguity. Therefore the actual finite zero set, on the
image under consideration, is exactly
$\chi_p(\mu_{2,p}^{-1}(0))$.

For $w=\chi_p(z)$, differentiate equivariance along the
compact action:
\[
       \mathrm i\delta_wS=d\chi_p(z)[\mathrm i\delta_zS].
\]
The derivative estimate and the positive bounded masses
imply $(1-C|z|)E_{p,z}\le E_{p,w}\le(1+C|z|)E_{p,z}$,
including exact equality of their kernels. Combine this
with Theorem~\ref{thm:complete-bundle-spectrum}.
The connection is the actual Chern connection of the
constructed holomorphic bundle, and its determinant and
nonzero Chern pairing are unchanged.
\end{proof}

The map $\chi_p$ is a real symplectic normalization.
It is not asserted to be holomorphic or to preserve
complex $G$-orbits. Its equivariance and injectivity do
give $K_{\chi_p(z)}=K_z$ as compact stabilizer groups.

\subsection{Balanced normal forms and stable factors}

Introduce the normalized arrows
\begin{equation}\label{eq:normalized-arrows}
 q=T_pz=(X,Y)=(\sqrt{c_+(p)}A,\sqrt{c_-(p)}B).
\end{equation}
The linear map $T_p$ is equivariant and uniformly
invertible, and the flat balanced equations are
\begin{equation}\label{eq:balanced-equations}
                    XX^*=Y^*Y,\qquad YY^*=X^*X.
\end{equation}
Singular-value decomposition makes $X=D$ diagonal
nonnegative by the $K$ action. Then $Y^*Y=YY^*=D^2$,
so $Y$ is normal and commutes with $D^2$. On a positive
singular-value block $d$, $Y/d$ is unitary and can be
diagonalized by the same unitary change at both vertices.
On the zero block it vanishes. Thus every balanced point
has a normal form
\begin{equation}\label{eq:balanced-normal-form}
 X=\diag(d_1,d_2),\qquad
 Y=\diag(d_1e^{\mathrm i\theta_1},d_2e^{\mathrm i\theta_2}),
                                      \qquad d_j\ge0.
\end{equation}
Let $m_\lambda$ count the entries with a given nonzero
$\lambda=d_j^2e^{\mathrm i\theta_j}$ and let $m_0$ count
the zero entries. The compact stabilizer is
\begin{equation}\label{eq:balanced-stabilizer}
 \prod_{\lambda\ne0}U(m_\lambda)\times U(m_0)\times U(m_0).
\end{equation}
Indeed the first stabilizer equation preserves $D^2$,
and the second separates the phase eigenspaces on each
positive block. The two zero vertex spaces transform
independently.

At the actual HE connection with coordinate $\chi_p(z)$,
Theorem~\ref{thm:complete-bundle-spectrum} identifies the
full Hermitian parallel algebra, after adding the scalar
identity, with the stabilizer algebra in
\eqref{eq:balanced-stabilizer}. Its complexification is
\begin{equation}\label{eq:parallel-matrix-algebra}
 \bigoplus_{\lambda\ne0}\operatorname{Mat}_{m_\lambda}(\CC)
       \oplus\operatorname{Mat}_{m_0}(\CC)
       \oplus\operatorname{Mat}_{m_0}(\CC).
\end{equation}
Every complex parallel section is the sum of Hermitian
real and imaginary parts, so this is the entire parallel
algebra. Its minimal orthogonal projections define actual
parallel holomorphic subbundles. A minimal projection in
a nonzero block has smooth bundle $L_p\oplus L_p^{-1}$
and only scalar parallel endomorphisms. The HE decomposition
theorem \cite[Theorem 5.8.3 and Proposition 5.7.6]{Kobayashi}
makes it slope stable, with stability understood against
coherent subsheaves. It has rank two, determinant $\OO$
and $c_2H=16$. Here the determinant is holomorphically
trivial because its integral first Chern class is zero
and $H^1(X_p,\OO)=0$.

Parallel matrix units identify the $m_\lambda$ copies.
Factors belonging to different nonzero blocks are
nonisomorphic within this bundle: a holomorphic
homomorphism of same-slope HE factors is parallel by
the section Bochner identity on their Hom bundle.
The zero blocks give $L_p$ and $L_p^{-1}$, each repeated
$m_0$ times; their holomorphic type is again fixed by
injectivity of the first Chern-class map on the Picard
group. Consequently the complete possibilities are
\begin{center}
\begin{tabular}{lrr}
\toprule
Normalized cycle type & Zero modes & Positive small modes\\
\midrule
Two distinct nonzero cycles & 1 & 6\\
One repeated nonzero cycle & 3 & 4\\
One nonzero cycle and one zero entry & 2 & 5\\
Both entries zero & 7 & 0\\
\bottomrule
\end{tabular}
\end{center}
Zero modes here are trace-free Hermitian endomorphisms.
Their number is $\sum_{\lambda\ne0}m_\lambda^2+2m_0^2-1$.
All remaining eigenvalues have the complementary gap
already proved. The table imposes no comparability
condition on its positive small eigenvalues.

\subsection{A correction within complex orbits}

We give the finite correction from \cite[Section 2]{W_Wall}
with the normalization and bounds needed here. For a finite
quiver of total vertex dimension $m$, put
\[
 (C_qU)_a=U_{h(a)}q_a-q_aU_{t(a)},\qquad
 \mathfrak a=\bigoplus_i\Herm(W_i),\qquad
 \omega_0(v,v')=-\operatorname{Im}\sum_a\tr(v_a^*v'_a).
\]
The incoming-minus-outgoing moment $\mu$ satisfies
\[
 \dd\tr(\mu U)=2\omega_0(\mathrm iC_qU,\cdot),
 \qquad\omega_0(\mathrm iC_qU,C_qU)=\|C_qU\|^2.
\]
Suppose $\Omega$ is a closed $K$-invariant form on a ball
$B_R$ and $\rho$ is its equivariant moment in this convention,
normalized at zero, with
\begin{equation}\label{eq:perturbed-moment-hypotheses}
 \rho(0)=0,\qquad
 \dd\tr(\rho U)=2\Omega(\mathrm iC_qU,\cdot),\qquad
                   \|\Omega-\omega_0\|_{B_R}\le\epsilon.
\end{equation}
No type-$(1,1)$ assumption on $\Omega$ is needed.

\begin{lemma}\label{lem:complex-orbit-correction}
Let $A_0=2m$ and
$\epsilon<\log2/(16A_0)$. If $x\in B_{R/4}$ satisfies
$\mu(x)=0$, there is a Hermitian tuple $S$ such that
$\rho(e^Sx)=0$ and $\|S\|_{\rm op}\le A_0\epsilon$.
It can be chosen in the centralizer of $K_x$.
The reverse conclusion holds starting from $\rho(y)=0$.
In a fixed complex orbit, either zero set contains at most
one compact orbit. These corrections induce inverse
continuous maps between the two local compact quotients
on nested neighbourhoods, including stabilizer changes.
\end{lemma}

\begin{proof}
For a tuple of orthogonal projections $P$, split $q=n+u$
into preserving and crossing arrow blocks. The circle
generated by $\mathrm iP$ fixes $n$. Both normalized
paired moments vanish on its fixed linear subspace, by
their Hamiltonian identities and their value at zero.
Integrate along the orthogonal segment $n+tu$ to obtain
\begin{equation}\label{eq:projection-relative-moment}
 |\tr((\rho-\mu)(q)P)|
 \le2\int_0^1t\epsilon\|u\|^2\,\dd t
                        =\epsilon\|C_qP\|^2.
\end{equation}
This estimate retains the crossing energy even when it
is arbitrarily small.

Let $H=K_x$, let
$\mathfrak h_x=\ker(C_x|_{\mathfrak a})$, and let
$\mathfrak c$ be the Hermitian centralizer of $H$.
Set $\mathfrak p=\mathfrak c\cap\mathfrak h_x^\perp$
for the trace pairing. On the fixed space $W^H$, both
moments take values in $\mathfrak p$: equivariance gives
$\mathfrak c$, while the generating fields for
$\mathfrak h_x$ vanish there, so their paired moments
vanish by normalization. All the homotopies below
therefore have values in this fixed real vector space.

For $S\in\mathfrak p$ of operator norm $a>0$, its total
trace is zero. Among its $m$ eigenvalues over all vertices
there is an adjacent gap at least $a/(m-1)$; the case
$m=1$ is trivial. Let $P$ project above such a gap.
Then $\|C_xP\|>0$. Otherwise $P,I-P\in\mathfrak h_x$,
as they also commute with $H$, and orthogonality gives
$\tr(SP)=\tr(S(I-P))=0$. The mean eigenvalues above
and below a strict gap cannot both be zero.

Take $a=A_0\epsilon$ and consider on the operator-norm
ball in $\mathfrak p$ the homotopy
\[
 F_t(S)=\mu(e^Sx)+t(\rho-\mu)(e^Sx),\qquad0\le t\le1.
\]
Every argument remains in $B_{R/2}$.
At a boundary point use the projection just chosen.
An upward crossing entry with eigenvalue difference
$d\ge a/(m-1)$ contributes $e^{2d}-1$ times its
squared norm to the paired moment difference; a downward
one contributes $1-e^{2d}$. Since $|d|\le2a$,
\begin{equation}\label{eq:degree-boundary-estimate}
 \begin{aligned}
 \tr((\mu(e^Sx)-\mu(x))P)
    &\ge2e^{-4a}\frac{a}{m-1}\|C_xP\|^2,\\
 |\tr((\rho-\mu)(e^Sx)P)|
    &\le\epsilon e^{4a}\|C_xP\|^2.
 \end{aligned}
\end{equation}
Our constants make the first coefficient larger, so
there is no boundary zero throughout the homotopy.
At $t=0$ the unique zero is $S=0$, because
\[
 \frac{\dd}{\dd t}\tr(\mu(e^{tS}x)S)
                         =2\|C_{e^{tS}x}S\|^2>0
\]
for nonzero $S\in\mathfrak p$. Its derivative at zero
has positive quadratic form $2\|C_xS\|^2$, hence degree
$+1$. Brouwer degree gives a zero at $t=1$.
If $\mathfrak p=0$ or $\epsilon=0$, the asserted zero
requires no correction.

For the reverse direction form the corresponding space
at $y$ and use
\[
 G_t(S)=\mu(e^Sy)-\mu(y)-t(\rho-\mu)(y).
\]
The initial degree calculation holds without assuming
$\mu(y)=0$. The first inequality in
\eqref{eq:degree-boundary-estimate} dominates
\eqref{eq:projection-relative-moment} at $y$.
At $t=1$ a zero satisfies $\mu(e^Sy)=\rho(y)=0$.

For uniqueness, polar-decompose a complex transformation
as $ke^U$ and use equivariance to compare $y$ and $e^Uy$.
The squared norm of $e^{tU}y$ is a sum of nonnegative
exponentials and is convex. If its endpoints are in
$B_R$, the entire path stays there. Along it,
\[
 \frac{\dd}{\dd t}\tr(\rho(e^{tU}y)U)
             \ge2(1-\epsilon)\|C_{e^{tU}y}U\|^2.
\]
If both endpoints are zeros, their paired values agree,
so $C_yU=0$ and they are compact-equivalent. The same
argument with $\epsilon=0$ proves flat uniqueness.

Finally a convergent sequence of starting representatives
has bounded correction tuples. A convergent subsequence
of those tuples produces a zero in the same complex orbit
as the limiting starting point. Uniqueness determines
its compact orbit independently of the subsequence.
This proves continuity even if the centralizer space
changes in the limit. With the corrections from $B_{R/4}$
to $B_{R/2}$, restriction makes the quotient maps inverse
on nested neighbourhoods. More explicitly the flat
quotient of $B_{R/8}$ maps onto its inverse image under
the reverse map on the perturbed quotient of $B_{R/4}$,
and this image contains the perturbed quotient of
$B_{R/16}$.
\end{proof}

In our coordinates \eqref{eq:normalized-arrows}, set
\begin{equation}\label{eq:orbit-correction-normalization}
 \rho_p(q)=V\Psi_p(T_p^{-1}q),\qquad
 \widetilde\Omega_p=-\tfrac12(T_p^{-1})^*\Omega_p.
\end{equation}
Equations \eqref{eq:actual-paired-hamiltonian} and
\eqref{eq:flat-pullback-form} give exactly
\eqref{eq:perturbed-moment-hypotheses}, with
$\widetilde\Omega_p(0)=\omega_0$ and
$\|\widetilde\Omega_p-\omega_0\|_{B_R}\le CR$.
Here $m=4$ and $A_0=8$.
Apply the lemma with $R=8|q_0|$, after choosing a common
outer domain. A flat-balanced $q_0$ has an actual zero
$q_1=e^Sq_0$ in the same complex orbit with
\begin{equation}\label{eq:small-orbit-correction}
 \|S\|_{\rm op}\le C|q_0|,\qquad
                    |q_1|\le e^{2C|q_0|}|q_0|.
\end{equation}
The origin has zero correction. The reverse assertion
holds with the same type of bound. These are distinct
maps from the real diffeomorphism $\chi_p$.

\subsection{Actual equivariance and semistability}

We record how these finite transformations act on the
holomorphic bundles. The effective étale roof of
Proposition~\ref{prop:effective-bundles} has a jointly
analytic group action near the identity and the zero
section. After two shrinks of the base and matrix
neighbourhood, compactness of the inner closure gives
a group neighbourhood $U_G$ of the identity such that
\begin{equation}\label{eq:local-actual-equivariance}
             \cE_{p,gw}\simeq\cE_{p,w}
\end{equation}
whenever $g\in U_G$ and the two representatives are in
the inner region. The lifts agree by uniqueness of the
inverse étale branch. If a piecewise smooth path $g(t)$
in $G$ stays in this coordinate region when applied to
$w$, subdivide its compact time interval so that successive
group increments lie in $U_G$, and compose these actual
isomorphisms. The resulting isomorphism can have arbitrarily
large norm. A path on $0<t\le1$ is treated on every
$[t,1]$ separately, without extending the group operation
to $t=0$.

The étale chart also identifies infinitesimal
automorphisms throughout a smaller neighbourhood:
\begin{equation}\label{eq:automorphism-dimensions}
 \operatorname{Lie}\Aut(\cE_{p,w})
             \simeq\operatorname{Lie}G_w,
 \qquad\dim_\CC\Aut(\cE_{p,w})=\dim_\CC G_w.
\end{equation}
Infinitesimal isomorphisms lift uniquely through the
étale diagonal; see \cite[Tag 0CJ1]{Stacks}.
Both groups are the open unit groups in finite-dimensional
complex endomorphism algebras, hence are smooth.
This proves the dimension assertion without presuming
that all isomorphisms between different chart points
come from $G$.

The algebraic family before pullback is locally free
and flat over its coefficient and bundle parameters.
Its central split bundle is Gieseker polystable for the
symmetric polarization. Indeed each line is stable as
a rank-one torsion-free sheaf on the integral central
fibre: a proper same-rank subsheaf subtracts a nonzero
torsion Hilbert polynomial, and a line has no nonzero
rank-zero subsheaf. The two reduced Hilbert polynomials
coincide by the Koszul calculation. Openness of Gieseker
semistability in a projective flat family
\cite[Proposition 2.3.1]{HL} therefore gives a product
neighbourhood of $(0,0)$ in which every bundle is
semistable. That source does not require smooth fibres.
Analytification and pullback to the prescribed period
ball preserve this open condition. On smooth fibres
Gieseker semistability implies slope semistability
\cite[Lemma 1.2.13]{HL}.

We shall also use that a slope-polystable,
Gieseker-semistable bundle is Gieseker polystable.
Its slope-stable summands are locally free and Gieseker
stable. Their reduced Hilbert polynomials are at most
that of the whole bundle, whereas their positive
rank-weighted average equals it. For all sufficiently
large integers the differences are nonnegative and their
weighted sum is zero. Each difference therefore vanishes
as a polynomial, which proves the assertion.

\begin{theorem}\label{thm:local-polystability}
There is a common matrix neighbourhood on the full small
smooth period base in which every $\cE_{p,w}$ is slope
semistable. It is slope polystable if and only if $w$ is
a semisimple two-arrow representation, equivalently its
affine complex orbit is closed. Every such object has an
actual HE representative in the family constructed above.
Its stable factors are the rank-two factors and line
factors already described, so no bundle in this rank-four
neighbourhood is slope stable. The chart contains an open
neighbourhood of the split object in the local bundle
stack, and the assertion applies to every polystable
object in this neighbourhood.
\end{theorem}

\begin{proof}
For a simple two-arrow representation $(X,Y)$, the
endomorphism $(XY,YX)$ commutes with both arrows and
hence is a scalar $\lambda$ by Schur's lemma. If
$\lambda\ne0$, the arrows are inverse isomorphisms up
to that scalar. Identifying the two spaces by one arrow
makes the other $\lambda I$, so simplicity forces
vertex dimensions $(1,1)$. If $\lambda=0$ and $X\ne0$,
then $(\im X,0)$ is a proper nonzero subrepresentation,
since $YX=0$; the same holds for $Y\ne0$. With both
arrows zero, a simple representation is one-dimensional
at just one vertex. Thus the simple types are
\[
 (1,0),\qquad(0,1),\qquad(1,1)\text{ with cycle }\lambda\ne0.
\]

Take any composition series and orthogonally split its
flags at the vertices. Each arrow sends block $j$ only
to blocks $i\le j$. Assign integer weight $\ell-j$ to
the $j$th block of a length-$\ell$ series.
The corresponding one-parameter group multiplies an
$i,j$ arrow block by $t^{j-i}$. Consequently
\begin{equation}\label{eq:bounded-semisimplification}
 q(t)=g(t)q\longrightarrow q_{\rm ss},\qquad
 |q(t)|^2=\sum_{i\le j}|t|^{2(j-i)}\|q_{ij}\|^2
                   \le|q|^2\quad(|t|\le1).
\end{equation}
The disc is polynomial, and its limit is the direct
sum of simple quotients. Orthogonal splitting ensures
the norm bound in the original metric.

Balance each nonzero $(1,1)$ pair $(x_j,y_j)$ in
$q_{\rm ss}$ by the positive ratio
$r_j=\sqrt{|y_j|/|x_j|}$ between vertex scalars.
Both resulting magnitudes are $\sqrt{|x_jy_j|}$, and
$2|x_jy_j|\le|x_j|^2+|y_j|^2$.
Along the Hermitian exponential path the squared norm
is convex, hence never exceeds its endpoint maximum.
After a unitary change the result has
\eqref{eq:balanced-normal-form}. If $q$ is semisimple,
this balanced point $q_0$ lies in its orbit and has
$|q_0|\le|q|$. To realize a bounded path directly from
$q$, choose $gq=q_0$ and write $g=ke^U$.
The exponential path stays in the ball by convexity,
and a compact-group path reaches $q_0$ at unchanged
norm. Equation \eqref{eq:local-actual-equivariance}
and the small correction
\eqref{eq:small-orbit-correction} give an actual HE
bundle isomorphic to $\cE_{p,T_p^{-1}q}$.
This proves polystability for semisimple parameters.

For arbitrary $q$, \eqref{eq:bounded-semisimplification}
pulls back the actual locally free family to a
holomorphic disc on the fixed smooth fibre. Each
nonzero member is isomorphic to the original bundle,
by the bounded-path realization on compact subintervals.
The central member is the semisimple limit, which has
an HE representative by the preceding argument.
Choose the initial neighbourhood with slack for the
final small correction, so all these parameters stay
in the common Gieseker-semistable region.

The projective moduli space of Gieseker-semistable sheaves
classifies S-equivalence
\cite[Theorems 4.3.3--4.3.4]{HL}. The analytified coarse
moduli map composed with this holomorphic disc is
constant away from zero and continuous into a separated
space, so it has the same value at zero. Thus the original
bundle and its semisimple limit are S-equivalent.
If the original bundle were slope polystable, both
would be Gieseker polystable by the preceding polynomial
argument. Uniqueness of the polystable representative
\cite[Proposition 1.5.2 and Definitions 1.5.3--1.5.4]{HL}
would make them isomorphic. If $q$ were nonsemisimple,
$q_{\rm ss}$ would lie in the boundary of its orbit.
An algebraic orbit is open in its irreducible closure;
every boundary orbit has smaller dimension. Hence
$\dim G_{q_{\rm ss}}>\dim G_q$, contradicting
\eqref{eq:automorphism-dimensions} for the isomorphic
bundles. Therefore polystability forces semisimplicity.
Closed affine orbits are exactly semisimple representations
by \cite[Proposition 3.2(i), $\theta=0$]{King}.

For the precise factor types in these complex-orbit
coordinates, choose the correction $S$ in the Hermitian
centralizer of $H=K_{q_0}$, as in
Lemma~\ref{lem:complex-orbit-correction}. Then
$H\subset K_{e^Sq_0}$. The complex stabilizer of the
normal form is
$\prod_{\lambda\ne0}\GL(m_\lambda,\CC)\times\GL(m_0,\CC)^2$,
whose complex dimension equals $\dim_\RR H$.
The stabilizer of $e^Sq_0$ is conjugate to it, and the
real dimension of a compact matrix subgroup cannot
exceed that complex dimension: its real Lie generators
are complex-linearly independent, since a Hermitian
matrix cannot equal a nonzero anti-Hermitian matrix.
The inclusion therefore gives equality of the compact
Lie algebras. By the exact kernel theorem, the actual
parallel algebra is precisely
\eqref{eq:parallel-matrix-algebra}. Its minimal projections
give the same stable rank-two factors and the two line
types, with the same multiplicities. Each rank-four
case has more than one stable summand. A stable bundle
would be polystable and hence in one of these cases,
which is impossible.

Finally the effective chart is étale as a morphism of
local stacks, and its atlas is smooth and open.
The image of a small analytic product neighbourhood
therefore contains an open neighbourhood of the split
object. Every polystable object in that image on a
smooth fibre has a representative in the chart; the
criterion just proved and the bounded correction give
its controlled HE representative. This is local
surjectivity of the actual deformation germ.
\end{proof}

The controlled representative is uniformly small.
A gauge taking an arbitrary unbalanced representative
to it can be arbitrarily ill-conditioned.
Injectivity on actual bundle isomorphism classes and
the harmonic metric on their quotient will be proved
in Section~\ref{sec:harmonic-quotient}; neither follows
from local surjectivity alone.
